% This is text associated with the open textbook "Open Electromagnetics"
% See immediately below for author, date
% This work is licensed under the Creative Commons Attribution-ShareAlike 4.0 International License. (CC BY-SA 4.0) 
% To view a copy of the license, visit http://creativecommons.org/licenses/by-sa/4.0/

%ORB TITLE Dielectric Breakdown
%ORB AUTHOR 0        # S.W. Ellingson, Virginia Tech (ellingson@vt.edu)
%ORB DATE 20180806   # yyyymmdd
%ORB ABSTRACT Dielectric Breakdown

\label{m0109_Dielectric_Breakdown} % <-- Must match name of this file and module directory name.  Don't mess with this!

The permittivity 
\index{permittivity}
of an ideal dielectric is independent of the magnitude of an applied electric field; the material is said to be ``linear.''\footnote{See Section~\ref{m0007_Properties_of_Materials} for a review of this concept.}
\index{linear (media)}
However, all practical dielectrics fail in this respect with sufficiently strong electric field.
Typically, the failure is abrupt and is observed as a sudden, dramatic increase in conductivity, signaling that electrons are being successfully dislodged from their host molecules.
The threshold value of the electric field intensity at which this occurs is known as the \emph{dielectric strength}, 
\index{dielectric strength}
and the sudden change in behavior observed in the presence of an electric field greater than this threshold value is known as \emph{dielectric breakdown}.
\index{dielectric breakdown}

Dielectric strength varies from about 3~MV/m for air to about 200~MV/m in mica 
\index{mica}
(a dielectric commonly used in capacitors).

Dielectric breakdown is typically accompanied by ``arcing,'' 
\index{arcing}
which is a sudden flow of current associated with the breakdown.
A well known example of this phenomenon is lightning, 
\index{lightning}
which occurs when charge is exchanged between sky and ground when air
\index{air}
(a dielectric) exhibits breakdown.
Dielectric breakdown in solids typically damages the material. 

{\bf Additional Reading:}
\begin{itemize}
\item ``\href{https://en.wikipedia.org/wiki/Electrical_breakdown}{Electrical Breakdown}'' on Wikipedia.
\end{itemize}


